\documentclass[11pt,a4paper]{article}
\usepackage[english,greek]{babel}
\usepackage[utf8]{inputenc}
\usepackage{nimbusserif}
\usepackage[T1]{fontenc}
\usepackage[left=1.50cm, right=1.50cm, top=2.00cm, bottom=2.00cm]{geometry}
\usepackage{amsmath}
\let\myBbbk\Bbbk
\let\Bbbk\relax
\usepackage[amsbb,subscriptcorrection,zswash,mtpcal,mtphrb,mtpfrak]{mtpro2}
\usepackage{graphicx,multicol,multirow,enumitem,tabularx,mathimatika,gensymb,venndiagram,hhline,longtable,tkz-euclide,fontawesome5,eurosym,tcolorbox,tabularray,diffcoeff}
\usepackage[explicit]{titlesec}
\tcbuselibrary{skins,theorems,breakable}
\newlist{rlist}{enumerate}{3}
\setlist[rlist]{itemsep=0mm,label=\roman*.}
\newlist{alist}{enumerate}{3}
\setlist[alist]{itemsep=0mm,label=\alph*.}
\newlist{balist}{enumerate}{3}
\setlist[balist]{itemsep=0mm,label=\bf\alph*.}
\newlist{Alist}{enumerate}{3}
\setlist[Alist]{itemsep=0mm,label=\Alph*.}
\newlist{bAlist}{enumerate}{3}
\setlist[bAlist]{itemsep=0mm,label=\bf\Alph*.}
\newlist{askhseis}{enumerate}{3}
\setlist[askhseis]{label={\Large\thesection}.\arabic*.}
\renewcommand{\textstigma}{\textsigma\texttau}
\newlist{thema}{enumerate}{3}
\setlist[thema]{label=\bf\large{ΘΕΜΑ \textcolor{black}{\Alph*}},itemsep=0mm,leftmargin=0cm,itemindent=18mm}
\newlist{erwthma}{enumerate}{3}
\setlist[erwthma]{label=\bf{\large{\textcolor{black}{\Alph{themai}.\arabic*}}},itemsep=0mm,leftmargin=0.8cm}

\newcommand{\kerkissans}[1]{{\fontfamily{maksf}\selectfont \textbf{#1}}}
\renewcommand{\textdexiakeraia}{}

\usepackage[
backend=biber,
style=alphabetic,
sorting=ynt
]{biblatex}

\DeclareTblrTemplate{caption}{nocaptemplate}{}
\DeclareTblrTemplate{capcont}{nocaptemplate}{}
\DeclareTblrTemplate{contfoot}{nocaptemplate}{}
\NewTblrTheme{mytabletheme}{
\SetTblrTemplate{caption}{nocaptemplate}{}
\SetTblrTemplate{capcont}{nocaptemplate}{}
\SetTblrTemplate{contfoot}{nocaptemplate}{}
}

\NewTblrEnviron{mytblr}
\SetTblrStyle{firsthead}{font=\bfseries}
\SetTblrStyle{firstfoot}{fg=red2}
\SetTblrOuter[mytblr]{theme=mytabletheme}
\SetTblrInner[mytblr]{
rowspec={t{7mm}},columns = {c},
width = 0.85\linewidth,
row{odd} = {bg=red9,fg=black,ht=8mm},
row{even} = {bg=red7,fg=black,ht=8mm},
hlines={white},vlines={white},
row{1} = {bg=red4, fg=white, font=\bfseries\fontfamily{maksf}},rowhead = 1,
hline{2} = {.7mm}, % midrule  
}
\newcounter{askhsh}
\setcounter{askhsh}{1}
\newcommand{\askhsh}{{\large\theaskhsh.}\ \addtocounter{askhsh}{1}}

\titleformat{\section}{\Large}{\kerkissans{\thesection}}{10pt}{\Large\kerkissans{#1}}

\setlength{\columnsep}{5mm}
\titleformat{\paragraph}
{\large}%
{}{0em}%
{\textcolor{red!80!black}{\faSquare\ \ \kerkissans{\bmath{#1}}}}
\setlength{\parindent}{0pt}
\titlespacing{\paragraph}{0mm}{2mm}{2mm}

\newcommand{\eng}[1]{\selectlanguage{english}#1\selectlanguage{greek}}

\begin{document}
\paragraph{Θέμα 5 - 2025}
Να επαληθευτεί το θεώρημα της απόκλισης 
\[\oiint_{S}{\vec{F}\cdot\hat{n}\d \sigma}=\iiint_{V}{\left(\vec{\nabla}\cdot\vec{F}\right)\d V}\] για το πεδίο $\vec{F}(x,y,z)=x\vec{i}+y\vec{j}+z\vec{k}$ όπου $S$ η επιφάνεια της σφαίρας $x^2+y^2+z^2=R^2$.\\\\
\textbf{ΛΥΣΗ}
\paragraph{Επιφανειακό ολοκλήρωμα}
\begin{enumerate}[label=\textbf{{Βήμα }\arabic*.},leftmargin=0mm,itemindent=14mm]
\item \textbf{Κάθετο διάνυσμα στην επιφάνεια}\\
Επειδή η επιφάνεια είναι σφαίρα μπορείς να πάρεις τον έτοιμο τύπο για το μοναδιαίο κάθετο διάνυσμα $\hat{n}$ που είναι παράλληλο με το διάνυσμα θέσης ενός σημείου της σφαίρας:
\[\hat{n}=\frac{(x,y,z)}{\sqrt{x^2+y^2+z^2}}=\frac{(x,y,z)}{R}\]
Αν θες να θυμάσαι μόνο το γενικό τύπο για το κάθετο διάνυσμα για οποιαδήποτε επιφάνεια θα είναι
\[\hat{n}=\frac{\vec{\nabla}S}{\left|\vec{\nabla}S\right|}=\frac{\left(\diffp{S}{x},\diffp{S}{y},\diffp{S}{z}\right)}{\sqrt{\left(\diffp{S}{x}\right)^2+\left(\diffp{S}{y}\right)^2}+\left(\diffp{S}{z}\right)^2}=\frac{(2x,2y,2z)}{\sqrt{4x^2+4y^2+4z^2}}=\frac{2(x,y,z)}{2\sqrt{x^2+y^2+y^2}}=\frac{(x,y,z)}{R}\]
\item \textbf{Εσωτερικό γινόμενο}\\
Έχουμε λοιπόν
\[\vec{F}\cdot\hat{n}=(x,y,z)\cdot\frac{(x,y,z)}{R}=\frac{x^2+y^2+z^2}{R}=\frac{R^2}{R}=R\]
\item \textbf{Υπολογισμός ολοκληρώματος}\\
\[\oiint_{S}{\vec{F}\cdot\hat{n}\d \sigma}=\oiint_{s}{R\d\sigma}=R\undercbrace{\oiint_{S}\d\sigma}_{\text{εμβαδόν σφαίρας}}=R\cdot4\pi R^2=4\pi R^3\]
\end{enumerate}
\paragraph{Τριπλό ολοκλήρωμα}
\begin{enumerate}[label=\textbf{{Βήμα }\arabic*.},leftmargin=0mm,itemindent=14mm]
\item \textbf{Εσωτερικό γινόμενο}\\
Έχουμε
\[ \vec{\nabla}\cdot\vec{F}=\diffp{\vec{F}}{x}+\diffp{\vec{F}}{y}+\diffp{\vec{F}}{z}=1+1+1=3 \]
\item \textbf{Υπολογισμός ολοκληρώματος}\\
\[\iiint_{V}{\left(\vec{\nabla}\cdot\vec{F}\right)\d V}=\iiint_{V}{3\d V}=3\undercbrace{\iiint_{V}{\d V}}_{\text{όγκος σφαίρας}}=3\cdot\frac{4}{3}\pi R^3=4\pi R^3\]
\end{enumerate}
Άρα επαληθεύεται.
\end{document}
